% Figure: tributary-area load path for the case-study floor beam.
% A row of joists at spacing s delivers a uniform line load to the
% interior beam over span L; the beam delivers half its load to each
% of two girders. The diagram shows the tributary width and the
% resulting line load.
\begin{figure}[htbp]
\centering
\begin{tikzpicture}[scale=0.9]
% plan view of a floor bay
% the interior beam (heavy) runs left-right
\draw[engblack, very thick] (0,0) -- (8,0);
\node[engblack, anchor=north, font=\scriptsize] at (4,-0.15) {interior beam, span $L$};
% girders at the two ends, running into the page (drawn vertical)
\draw[engblue, line width=3pt] (0,-1.6) -- (0,1.6);
\draw[engblue, line width=3pt] (8,-1.6) -- (8,1.6);
\node[engblue, anchor=east, font=\scriptsize] at (-0.1,1.4) {girder};
\node[engblue, anchor=west, font=\scriptsize] at (8.1,1.4) {girder};
% joists running parallel to the girders, spaced along the beam
\foreach \x in {1,2,3,4,5,6,7}{
  \draw[enggray, thick] (\x,-1.3) -- (\x,1.3);
}
\node[enggray, anchor=south, font=\scriptsize] at (3.5,1.3) {joists at spacing $s$};
% tributary band of the interior beam (half a joist span each side)
\fill[engamber!18] (0,-0.8) rectangle (8,0.8);
\draw[<->, engamber] (8.6,-0.8) -- (8.6,0.8);
\node[engamber, anchor=west, font=\scriptsize] at (8.7,0) {tributary width $w_t$};
% resultant line load arrows on the beam
\foreach \x in {0.5,1.5,2.5,3.5,4.5,5.5,6.5,7.5}{
  \draw[engred, ->] (\x,1.1) -- (\x,0.12);
}
\node[engred, anchor=south, font=\scriptsize] at (4,1.15) {line load $w_0 = q\,w_t$};
\end{tikzpicture}
\caption[Tributary load path for a floor beam]{Plan of one floor bay.
The floor pressure $q$ is collected by the joists and delivered to the
interior beam over a tributary width $w_t$ equal to half the joist span
on each side of the beam. The beam therefore carries a uniform line
load $w_0 = q\,w_t$, and passes half of its total load to each of the
two girders at its ends. The first step of every floor-beam design is
this load path: pressure to tributary width to line load to end
reactions, before any stress or deflection is computed.}
\label{fig:vol03ch08:floor-beam-loadpath}
\end{figure}
